\chapter{Mediation Analysis with One Mediator}\label{chapter:sample}

We begin by formulating the causal mediation problem in terms of the potential outcomes
framework~\citep{rubin1974}. We first review the potential outcomes framework for "simple"
causal inference problems with no mediation, where the goal is to infer the causal 
effect of a particular change in a treatment variable on an outcome variable. We then 
extend this framework to the case of mediation analysis with one mediator and discuss the 
various mediational effects that have been proposed in the literature. Throughout this section, we use
the running example from~\cite{morenobetancur2018} introduced in Section~\ref{chapter:intro} to demonstrate
the different types of effects we are interested in estimating. 

\section{Potential Outcomes}

Suppose we are weighing whether a certain anti-depressant that has been developed is effective in treating
depression. To evaluate this question, we have access to observational data from a study on $N$
patients diagnosed with clinical depression indexed by $i = 1, 2, ..., N$. It is imagined that
these patients were a random sample taken from a larger population of clinically despressed patients. Each patient,
at the start of the study period, was given the option of beginning uptake of the anti-depressant. 
We have the following data:
\begin{itemize}[noitemsep,topsep=0pt]
\item The treatment variable $A_{i}$ denotes a binary treatment variable that equals $1$ if the 
unit takes a certain dose of the anti-depressant and equals $0$ otherwise. 
\item The outcome variable $Y_{i}$ represents the patient's score on a questionnaire evaluating the severity 
of depression in the patient that is filled out months after the patient begins or does not begin treatment. 
\item Covariates $C_{i}$ are unit-specific characteristics (e.g. sex, medical history) that are potentially 
related to the patient's propensity to take treatment and the strength of their depression 6 months later. 
\end{itemize}

Let $\mathcal{A}$, $\mathcal{Y}$, and $\mathcal{C}$ denote the supports of $A_{i}$, $Y_{i}$, and $C_{i}$,
respectively. 

Under the potential outcomes framework, it is imagined that each unit in our sample also has
two potential values of $Y_{i}$ --- their realized outcome under the treatment level they actually received
and the outcome they would have experienced had they received the other treatment level. Let $Y_{i}(1)$
and $Y_{i}(0)$ denote the outcome values unit $i$ would obtain under each treatment level. One of these
potential outcomes will always be observed and will be equal to the realized $Y_{i}$, while the other
will always be missing. 

The \textit{unit level treatment effect} of $A$ on $Y$ is equal to $Y_{i}(1) - Y_{i}(0)$ and represents 
the effect of taking anti-depressants on unit $i$'s severity of depression. This
effect is never identified for individual units because one potential outcome for each unit is always
unobserved ---~\cite{rubin1974} and~\cite{holland1986} call this "the fundamental problem of causal 
inference". However, different population-level causal estimands are identified under certain assumptions. 
One such estimand we focus on is the \textit{average treatment effect} (ATE), given by 
$\mathbb{E}[Y_{i}(a^{*}) - Y_{i}(a)]$.

The ATE represents the average causal effect of the treatment on the outcome in the larger population. In
our example, this represents the average difference of patients' severity of depression in worlds where they 
are induced to take the anti-depressant medication compared to worlds in which they 
do not take the medication, holding all other relevant factors constant. The ATE, and other 
important estimands, can be identified given the following two assumptions~\citep{imbens2015}: 

\begin{assumption}[Stable Unit Treatment Value Assumption]\label{assm:SUTVA}
\indent
\begin{enumerate}[noitemsep,topsep=0pt]
\item \textbf{Consistency:} $Y_{i}(a) = Y_{i}$ if unit $i$ receives treatment level $a \in \mathcal{A}$. 
\item \textbf{No interference:} $Y_{i}(a)$ depends only on the treatment level $a$ for unit $i$; it does not depend on 
the treatment level for any other unit $j \neq i$. 
\end{enumerate}
\end{assumption}

\begin{assumption}[Conditional Ignorability]\label{assm:ignorability} After controlling for some set of 
observable covariates $C$, treatment level is independent of unit potential 
outcomes: $Y_{i}(a) \indep A_{i} | C_{i}$ for all $a \in \mathcal{A}$. 
\end{assumption}

The Stable Unit Treatment Value Assumption (SUTVA)~\citep{rubin1980a} gives technical conditions under which
the ATE can be identified. It requires that there is no interference between units, meaning that a unit's 
treatment status does not affect the potential outcomes of other units and vice versa. In this example, a violation
might occur if we believed depression within a household was "contagious", meaning that whether the people a patient
was living with received treatment or not (if they were included in the study) affected the patient's outcome.  
Additionally, there are to be no "hidden versions" of treatments, meaning that a given treatment 
level cannot map to multiple different treatments for a unit. For example, in the birth control example from 
Section~\ref{chapter:intro}, a violation would occur if some women who are recording as receiving the treatment take
expired birth control pills that have lower efficacy. Violations can also occur in experimental settings where differences
in the mode of treatment assignment itself --- for example, the attitude of the doctor or lab assistant administering 
treatment --- might affect potential outcomes.

The extension of the SUTVA assumption to mediation analysis is not to be made lightly. In particular,~\cite{imaietal2013} 
discusses the difficulty of satisfying the consistency assumption in hypothetical 
experimental designs for identifying natural direct and indirect effects. However, violations of SUTVA are not 
the focus of this dissertation, and SUTVA and its extensions will be assumed without further comment 
in the rest of this dissertation. 

The conditional ignorability assumption~\citep{rubin1978} allows us to identify estimands like the ATE because while both 
potential outcomes for a single unit can never be observed, ignorability allows us to treat these missing potential 
outcomes as if they were "missing at random" because there will be no bias driving which potential outcomes are revealed
versus hidden~\citep{imbens2015} after controlling for confounders $C_{i}$. This assumption
rules out selection bias: for example, it rules out the possibility that patients with an unobserved lower desire to 
make positive lifestyle changes are also less likely to take the anti-depressants. Alternatively, a randomized experiment
where all study participants were randomly assigned to either take the anti-depressant or a placebo would also satisfy 
this assumption. Thus, the difference in a covariate-weighted average of the realized potential outcomes at both
treatment levels is an unbiased estimator of the ATE.

Throughout this paper, we will also present directed acyclic graph (DAG) representations of the causal structures 
we are developing within the potential outcomes framework. The non-parametric structural equation 
modeling or DAG framework of~\cite{pearl2000} is an alternative method 
for approaching causal inference problems that has been adapted by many applied researchers, especially in 
epidemiology~\citep{hernan2024}, and is particularly helpful for visually encoding the types of identification 
assumptions made above. Much of the early work on mediation analysis (e.g.~\citep{pearl2001,avinetal2005}) was
conducted within the DAG framework. We find DAG representations to be helpful in gaining intuition for the types of 
complicated causal structures with many variables that are encountered in mediation analysis, but these graphs are 
presented as illustrative aids only. 

\begin{figure}[ht]
\centering
\begin{tikzpicture}

  % Define nodes
  \node[obs]                               (A) {$\mathbf{A}$};
  \node[obs, right=1.5cm of A]  (Y) {$\mathbf{Y}$};
  \node[obs, left=1.5cm of A]  (C) {$\mathbf{C}$};

  % Connect the nodes
  \edge {A} {Y} ; %
  \draw[->] (C) to[out=315, in=225] (Y); %
  \edge {C} {A} ; %

\end{tikzpicture}
\caption[]{DAG encoding the conditional ignorability assumption}
\label{fig:one_mediator_dag}
\end{figure}

% Potentially add discussion of the difference between randomized experiments and observational studies here; can
% be used to emphasize the difficulty of justifying the assumptions for natural effect identification


\section{Natural Effects}

\subsection{Motivation}

Suppose we now observe a small effect of the anti-depressant on depression and are worried 
that a deleterious side effect associated with the anti-depressant, chronic headaches, might be interfering 
with the positive effects of the drug. In this situation, there are two possibilities: 

\begin{itemize}[noitemsep,topsep=0pt]
	\item The drug is effective at treating existing depressive symptoms, but the chronic pain from the side-effects 
	cancel out this benefit and provide a new mental health challenge. 
	\item The headaches are a minor mental health issue, but the drug is not effective at treating the original 
	symptoms. 
\end{itemize}

While the observed causal effect of the anti-depressant is the same in both situations, which of the two situations is 
operative has important implications for evaluating how promising the drug is. If we are in the first world,
a drug manufacturer might consider investing resources in reducing the side effects of an otherwise promising drug;
if we are in the second world, the anti-depressant does not seem like a promising treatment to invest in. 


\subsection{Model}

We now define a mediator variable $M_{i}$, which represents a continuous measure of whether the intensity of chronic
headache pain a patient experienced during the study period. Let $\mathcal{M}$ represent the support of $M$. 
We also define potential outcomes (the definition of which 
we now extend to include potential values for both mediators and outcomes) over both $M_{i}$ and
$Y_{i}$:
\begin{itemize}[noitemsep,topsep=0pt]
	\item $M_{i}(0)$ --- The level of headache a patient would experience if they did not take anti-depressants.
	\item $M_{i}(1)$ --- The level of headache a patient would experience if they took anti-depressants.
	\item $Y_{i}(0,M_{i}(0))$ --- Depressive symptoms a patient would experience if they were untreated and had 
	the level of headache they would experience if they were not treated. 
	\item $Y_{i}(1,M_{i}(1))$ --- Depressive symptoms a patient would experience if they were treated and had 
	the level of headache they would experience if they were treated. 
	\item $Y_{i}(1,M_{i}(0))$ --- Depressive symptoms a patient would experience if they were treated but had 
	the level of headache they would experience if they were untreated.
	\item $Y_{i}(0,M_{i}(1))$ --- Depressive symptoms a patient would experience if they were untreated but had 
	the level of headache they would experience if they were treated.
\end{itemize}

If a unit does not receive treatment, we observe their realized outcomes $M_{i}(0)$ and $Y_{i}(0,M_{i}(0))$; if they
are treated, we observe $M_{i}(1)$ and $Y_{i}(0,M_{i}(1))$. Note that we can never observe $Y_{i}(1,M_{i}(0))$ or
$Y_{i}(0,M_{i}(1))$ because a unit can only be assigned a treatment level once, but evaluating $Y_{i}(t',M_{i}(t))$ for 
$t' \neq t$ would require knowing both the unit's outcome under treatment level $t'$ and their mediator value under
treatment level $t$~\citep{robinsgreenland1992,pearl2001}. Having defined this expanded set of potential outcomes, we
now define a set of estimands that can help define the path-specific causal effects we are interested in:
\begin{itemize}[noitemsep,topsep=0pt]
	\item $\mathbb{E}[Y_{i}(1,M_{i}(0)) - Y_{i}(0,M_{i}(0))]$ --- The \textit{pure direct effect} (PDE)
	is the average causal effect that receiving treatment has on the outcome, keeping the level of the mediator constant 
	for each unit at the level of the mediator that would have been realized had the unit not received treatment. 
	\item $\mathbb{E}[Y_{i}(0,M_{i}(1)) - Y_{i}(0,M_{i}(0))]$ --- The \textit{pure indirect effect} (PIE)
	is the average causal effect that changes in the mediator that are induced from changes in the treatment level have
	on the outcome, evaluated at the untreated level of the treatment.  
	\item $\mathbb{E}[Y_{i}(1,M_{i}(1)) - Y_{i}(1,M_{i}(0))]$ --- The \textit{total indirect effect} (TIE)
	is the average causal effect that changes in the mediator that are induced from changes in the treatment level have
	on the outcome, evaluated at the treated level of the treatment.  
	\item $\mathbb{E}[Y_{i}(1,m) - Y_{i}(0,m)]$ --- The \textit{controlled direct effect} (CDE) at $M=m$ is the 
	average causal effect that receiving treatment has on the outcome, keeping the level of the mediator fixed at a 
	particular level $m$ for all units. There is a (potentially different) CDE defined at all possible values of the 
	mediator. 
	\item $\mathbb{E}[Y_{i}(1,M_{i}(1)) - Y_{i}(0,M_{i}(0))]$ --- The \textit{total effect} (TE) is 
	the average causal effect of a change in the treatment on the outcome. This is equal to the ATE and is identified
	under the same assumptions. 

\end{itemize}
% Pure vs. total indirect comes from Robins and Greenland 1992
% Robins and Greenland 1992 interaction example (can drop to binary mediator/outcome if helpful)
% Potential table of attribution for each effect
% Vanderweele decomposition


% Emphasize that these are rough definitions
% Main sources for the development of these effects: Robins and Greenland (1992); Pearl (2001)
% Imai, Keele, Yamamoto 2010: sequential ignorability assumption (read this paper again closely)
% Imai, Tingley, Yamamoto 2013 on experimental design? Sensitivity analysis?
% Various citations for issues with natural effects: Vansteelandt and Vanderweele 2009 (interpretation/un-testability), 
% Vanderweele et al 2014 and Avin et al (exposure-induced confounding) 

% Section 4: look at Robins and Greenland 1992, do a similar breakdown of the types of patients in a binary setting 
% in the two ordered mediators case; generalize, using Vanderweele 2013/2014, to the interventional effects case

% Should the Vanderweele decomposition framework enter the picture in Section 3? This would allow me to use the framework
% when discussing natural effects from the beginning, but it also would require a serious re-write as well as an 
% acknowledgement that some of the effects definitions (which ones?) don't make much sense (except that they do!)

% I should do this, but how to frame?

% Structure:
% - Potential outcomes framework
% - Mediation analysis with one mediator 
%.     - Notation, causal story (informal direct and indirect effects) (Pearl 2001)
%.     - Natural effects (pure vs. indirect )
%             - Breakdown of types of patients in the applied example (Robins and Greenland 1992)
%             - Mapping these types of patients to pure vs. natural direct/indirect effect 
%               definitions (Robins and Greenland 1992; Pearl 2001)
%             - Handling interaction terms (Vanderweele 2013/2014)
%             - Assumptions: standard ones + sequential ignorability, causal graph (Imai et al 2010)
%             - Problems: interpretational (Vanderweele and Vansteelandt 2009; briefly Robins and Greenland 2003) 
%               and mathematical (cross-world independence; exposure induced confounding) 
%               (Imai et al 2010; Avin et al 2005; Vanderweele et al 2014)
%.            - Mention/define controlled direct effects somewhere; this can be part of the 4-way decomposition from
%               Vanderweele 2014
%.            - Emphasize throughout (citing Robins and Greenland, Vanderweele stuff) the importance of the substantive
%               properties of the problem; this will pre-emptively justify the approach taken in Section 4
%.     - Interventional effects with one mediator 
%
% I should incorporate an applied example (probably medical is easiest) throughout my definition of the effects; will make
% it easiest to break down the sub-types of effects like in Robins and Greenland and the Vanderweele papers
% Social mobility example? I don't think this is the best because of the lack of credibility of a causal story. Maybe an
% epi example? Moreno-Betancur headache/anti-depressant causal story seems the best, because it can generalize to Section 4
%
% Differences between Section 3 and Section 4: Section 3 is a pure interpretation of the existing literature while
% Section 4 attempts to advance a "novel" framework that can clarify/make commensurable the various 
% multiple mediator effect definitions 



\begin{figure}[ht]
\centering
\begin{tikzpicture}

  % Define nodes
  \node[obs]                               (A) {$\mathbf{A}$};
  \node[obs, right=of A, yshift=1.5cm] (M) {$\mathbf{M}$};
  \node[obs, right=2.5cm of A]  (Y) {$\mathbf{Y}$};
  \node[obs, left=of A]  (C) {$\mathbf{C}$};

  % Connect the nodes
  \edge {A,M} {Y} ; %
  \draw[->] (C) to[out=315, in=225] (Y); %
  \edge {C,A} {M} ; %
  \edge {C} {A} ; %

\end{tikzpicture}
\caption[]{Mediation analysis with one mediator}
\label{fig:one_mediator_dag}
\end{figure}
